\documentclass[aspectratio=169,11pt,t]{beamer}
\input{../../../theme/ksa-cs-v2.tex}
\usepackage{fontawesome5}
\newcommand{\paperlink}[1]{\href{#1}{\textcolor{ksablue}{\faLink}}}
% 원본 덱 (ksa-theme) 의 그림·실습은 이 폴더로 합쳤다 (옛 원본: slides/old/v2/)
\setbeamerfont{author}{size=\normalsize}
\setbeamerfont{institute}{size=\small}
\graphicspath{{figs/}{../../../assets/}}
\usetikzlibrary{tikzmark,shapes.callouts,positioning}
\tikzset{bub/.style={draw=ksagolddk,line width=1pt,fill=ksagolddk!10,rounded corners=3pt,inner sep=4pt,font=\normalsize,align=left},
         bubarr/.style={-latex,ksagolddk,line width=1.2pt}}
\definecolor{fx}{HTML}{D62728}\definecolor{fsh}{HTML}{2E3192}\definecolor{fbeta}{HTML}{7A6A3F}\definecolor{fdist}{HTML}{000000}\definecolor{fbd}{HTML}{555555}
\makeatletter
\newcommand{\markpage}[1]{\protected@write\@auxout{}{\string\newlabel{#1}{{\the\c@framenumber}{}}}\hypertarget{#1}{}}
\newcommand{\pref}[1]{\hyperlink{#1}{\textcolor{ksagolddk}{(p.\ref{#1})}}}
\makeatother

\renewcommand{\ksacredit}{Machine Learning and Deep Learning \,\textperiodcentered\, 서포트 벡터 머신}
\title{SVM 1: 최대 마진 (Support Vector Machine)}
\subtitle{초평면 · 최대 마진 · 학습}
\author{Machine Learning and Deep Learning}
\institute{KSA 2027-Spring}
\date{Sumin Han (suminhan@ksa.hs.kr)}
\begin{document}
{\setbeamertemplate{footline}{}\begin{frame}[c,noframenumbering]\titlepage\end{frame}}

\section{질문}
% --- 오늘의 질문
\begin{frame}{오늘의 질문}
\vfill
\begin{columns}[c]
\begin{column}{0.44\textwidth}\centering\includegraphics[width=0.95\linewidth]{sv_hook.png}\end{column}
\begin{column}{0.52\textwidth}
{\large 세 선 모두 \textbf{완벽하게} 가른다}\\[0.8em]
{\Large \kq{어느 선이 가장 좋을까?}}
\end{column}
\end{columns}
\vfill
\note{빨간 점과 파란 점이 있습니다. 둘을 가르는 직선을 세 개 그었는데, 세 선 모두 학습 데이터를 하나도 틀리지 않고 완벽하게 가릅니다. 학습 오차로는 셋 다 0입니다. 그런데 어느 선이 가장 좋을까요? 직관적으로는 가운데 검은 선이 가장 안심이 됩니다. 왜 그런지, 그리고 그 선을 어떻게 찾는지가 오늘의 주제입니다.}
\end{frame}

% --- 데이터는 먼저 z-score
\begin{frame}{먼저: 점들은 z-score normalization 해 둔다}
\vfill
\centering\includegraphics[width=0.8\linewidth]{sv_raw_z.png}\par\vskip0.2em
{\large $z = (x - \text{평균}) / \text{표준편차}$ \quad (축마다) \quad 이 덱의 점은 모두 \kb{z-score 한 값}}
\vfill
\note{시작하기 전에 데이터를 한 번 손봅니다. 왼쪽이 원래 값입니다. 각 축에서 평균을 빼고 표준편차로 나누면 오른쪽처럼 점들이 원점 주변에 크기 1 정도로 모입니다. 이것을 z-score라고 합니다. x1은 평균 5, 표준편차 약 1이고, x2는 평균 5, 표준편차 약 3입니다. 세로로 3배 더 퍼져 있던 것이 z-score 후에는 두 축 모두 크기 1로 맞춰집니다. 오늘 덱의 점 20개, 그림, 숫자는 모두 이 z-score 한 값입니다. 소수 첫째 자리로 반올림했습니다. 모양은 그대로이고 눈금만 바뀝니다. 왜 이렇게 하는지는 학습 부분에서 봅니다. 하지 않으면 학습이 아예 실패할 수 있습니다. sklearn으로 SVM을 쓸 때도 StandardScaler로 z-score를 먼저 하는 것이 기본입니다.}
\end{frame}

% --- 파랑 뒤집기 → 평균 (한 쪽에 세 장)
\begin{frame}{가장 쉬운 방법: 점에 $y_k$ 를 곱해 평균 내기}
\markpage{meanvec}
\vfill
\centering\includegraphics[width=\linewidth]{sv_flip_mean.png}\par\vskip0.2em
{\large 파랑은 \kb{원점 기준 180° 회전} \;$\to$\; $\mathbf{w} = \dfrac{1}{p}\sum_{k=1}^{p} y_k\mathbf{x}_k = (0.83,\ 0.91)$, \;$b = 0$}
\vfill
\note{가장 먼저 떠오르는 방법부터 해 봅시다. 세 그림을 왼쪽부터 차례로 봅니다. 첫째 그림은 원래 점 20개입니다. 둘째 그림, 파란 점의 좌표 두 개에 모두 마이너스를 붙입니다. 라벨 y가 파랑은 마이너스 1이니, y k 곱하기 x k입니다. 좌표 두 개의 부호를 모두 바꾸면 원점을 중심으로 180도 돌린 것, 점대칭입니다. 선을 거울로 비추는 대칭이 아닙니다. 예를 들어 파란 점 마이너스 1, 마이너스 0.6은 1, 0.6으로 갑니다. 빈 네모가 뒤집은 파란 점이고, 이제 20개가 모두 빨강 쪽, 한쪽에 모입니다. 셋째 그림, 원점에서 각 점까지의 화살표, 즉 벡터 20개의 평균을 내면 금색 화살표, 0.83, 0.91입니다. 이 화살표를 w로 쓰고, w에 수직이고 원점을 지나는 선을 긋습니다. 점들을 z-score 해서 가운데가 원점이니 b는 0으로 둡니다. 계산은 평균 한 번이면 끝납니다. 뒤에서 나오는 거울로 접기는 경계를 거울 삼는 것이라 다른 이야기입니다. 그리고 뒤에서 보겠지만 이 선은 SVM 학습의 1걸음 뒤 첫 선과 똑같습니다.}
\end{frame}

% --- 다듬기: 평균 선 vs SVM
\begin{frame}{SVM: 가장 가까운 점 기준으로 다듬는다}
\vfill
\centering\includegraphics[width=0.78\linewidth]{sv_centroid_vs.png}\par\vskip0.2em
{\large 평균 선: 빨강 쪽으로 \kq{치우침} (0.62 vs 0.70) \quad SVM: 양쪽 \kb{0.67}}
\vfill
\note{그런데 이 평균 선에서 가장 가까운 점까지 거리를 재 보면, 빨간 점 0.7, 0.2까지는 0.62, 파란 점 마이너스 0.6, 마이너스 0.4까지는 0.70입니다. 평균은 먼 점까지 모두 섞어서 정하니, 경계 가까이 있는 점 기준으로 보면 한쪽으로 치우쳐 있습니다. SVM은 평균 대신 경계에 가장 가까운 점들만 보고, 그 사이의 한가운데로 선을 옮깁니다. 그러면 양쪽 모두 0.67로 같아지고, 이것이 가장 가까운 점까지 거리를 가장 크게 만든 선입니다. 회색 점선이 평균 선, 검은 선이 SVM입니다. 뒤에서 볼 학습은 0에서 출발하면 첫 걸음이 바로 이 평균 선이 되고, 그 뒤 가장 가까운 점들이 밀면서 SVM 선으로 다듬어 갑니다.}
\end{frame}

% --- 로지스틱회귀 vs SVM
\begin{frame}{로지스틱회귀 vs SVM}
\vfill
\centering
\renewcommand{\arraystretch}{2.2}
{\LARGE\begin{tabular}{l|l}
로지스틱회귀 & \textbf{모든 점}이 선을 민다 (틀릴수록 세게)\\ \hline
\kb{SVM} & \kb{가장 가까운 점}만 선을 민다
\end{tabular}}\\[1.2em]
{\large 값 계산은 둘 다 \textbf{모든 점}에 대해 한다}
\vfill
\note{세 방법을 학습 한 걸음으로 비교합니다. 세 방법 모두 한 걸음은 같은 꼴입니다. w에 학습률 곱하기 색 곱하기 좌표를 더하는데, 점마다 그 점이 미는 세기를 곱합니다. 다른 것은 이 세기뿐입니다. 모든 점의 평균 선은 모든 점이 똑같이 1로 밉니다. 먼 점도 가까운 점도 같은 무게입니다. 로지스틱회귀는 시그모이드로 계산한 틀릴 확률만큼 밉니다. 경계에서 멀고 맞는 점은 틀릴 확률이 0에 가까워 거의 안 밀고, 경계 근처는 0.5쯤, 틀린 점은 1에 가깝게 세게 밉니다. 다만 먼 점도 0은 아니라서 아주 조금씩은 계속 밉니다. 로지스틱 손실, 크로스 엔트로피를 미분하면 바로 이 꼴이 나옵니다. SVM도 가장 가까운 점을 찾으려면 매 걸음 20개 점의 값을 모두 계산해야 합니다. 계산은 모두 하고, 그중 값이 가장 작은 점, 즉 경계에 가장 가까운 점만 밀고 나머지는 미는 세기가 정확히 0입니다. 그래서 학습이 끝나면 가장 가까운 점 몇 개, 서포트 벡터만 남습니다. 시험으로 치면 로지스틱은 이미 잘하는 학생은 거의 안 세고 못하는 학생일수록 크게 세는 평균, SVM은 그 거의 안 셈을 아예 안 셈으로 끊어 꼴찌 몇 명만 보는 것입니다.}
\end{frame}

% --- 원 논문
\begin{frame}{SVM 의 원 논문}
\vfill
\centering
\renewcommand{\arraystretch}{1.9}
{\large\begin{tabular}{c@{\hspace{1em}}p{5.3cm}l}
\textbf{1992} & Boser, Guyon, Vapnik & \kb{최대 마진} 학습 + \kb{커널}\\[-0.4em]
 & \multicolumn{2}{l}{\small \textit{A training algorithm for optimal margin classifiers} (COLT '92) \paperlink{https://doi.org/10.1145/130385.130401}}\\[0.6em]
\textbf{1995} & Cortes, Vapnik & \kb{소프트 마진}, ``Support Vector'' 라는 이름\\[-0.4em]
 & \multicolumn{2}{l}{\small \textit{Support-vector networks} (\textit{Machine Learning} 20(3)) \paperlink{https://doi.org/10.1007/BF00994018}}
\end{tabular}}\\[1.2em]
{\large 덱의 기호 ($D(\mathbf{x})$, $M$, $y_k$ · 다음 시간 $\alpha_k$, $\xi_k$, $C$) 는 \kb{두 논문 그대로}}
\vfill
\note{오늘 배울 서포트 벡터 머신이 어디서 왔는지부터 봅시다. 원 논문은 보통 두 편을 같이 꼽습니다. 첫째는 1992년, AT&T 벨 연구소의 Guyon과 Vapnik, 그리고 Boser가 COLT라는 학회에 낸 A training algorithm for optimal margin classifiers입니다. 경계와 가장 가까운 점 사이의 여백, 마진을 가장 크게 하는 학습법을 내고, 커널을 써서 직선이 아닌 휘어진 경계까지 같은 방법으로 학습할 수 있다는 것을 보였습니다. 오늘 배우는 최대 마진과 세 번째 시간의 커널이 이 논문입니다. 최대 마진으로 직선을 고르는 생각 자체는 1960년대 Vapnik의 generalized portrait 방법까지 거슬러 올라가고, 커널로 점곱을 바꾸는 생각도 1964년 Aizerman 등의 논문에 있었습니다. 이 논문은 둘을 합쳤습니다. 둘째는 1995년 Cortes와 Vapnik이 Machine Learning 저널에 낸 Support-vector networks입니다. 1992년 논문은 두 무리가 직선으로 깔끔하게 갈릴 때만 다뤘는데, 이 논문이 몇 점은 넘어가도 되는 소프트 마진을 더해 실제 데이터에 쓸 수 있게 했고, 서포트 벡터라는 이름도 이 논문 제목에서 나왔습니다. 두 논문 모두 손글씨 숫자 인식으로 실험했습니다. 이 덱의 기호는 두 논문을 그대로 따릅니다. 결정 함수 D, 마진 대문자 M, 점 번호 k와 라벨 y k는 1992년 논문, 슬랙 크사이와 벌점 C는 1995년 논문 기호입니다.}
\end{frame}

% --- 비하인드
\begin{frame}{비하인드: AT\&T 벨 연구소, 1990년대 초}
\vfill
\centering
\renewcommand{\arraystretch}{1.45}
{\large\begin{tabular}{rl}
\textbf{커널은 누구 생각?} & Vapnik: ``특징을 곱해서 늘리자'' $\to$ Guyon: ``\kb{커널}로 하자''\\[0.8em]
\textbf{1995년 내기} & 신경망은 잘 맞히는데 \kq{왜 되는지 이론이 없던} 때\\
 & Jackel: ``2000년까지 \kb{이유를 알게 된다}''\\[-0.3em]
 & \hspace*{1.5em}{\normalsize $\to$ \kq{틀림}: 지금도 다 못 풀었다}\\
 & Vapnik: ``2005년엔 신경망이 \kb{안 쓰인다}''\\[-0.3em]
 & \hspace*{1.5em}{\normalsize $\to$ \kq{틀림}: 지금 가장 많이 쓰는 모델이 됐다}
\end{tabular}}\\[1.2em]
{\large 결과: \kq{둘 다 졌다}}
\vfill
\note{비하인드 이야기 두 개입니다. SVM은 1990년대 초 미국 AT&T 벨 연구소에서 나왔습니다. Vapnik은 소련에서 1960년대부터 학습 이론을 연구하다가 1990년에 미국으로 건너와 이 연구소에 합류했습니다. 첫째, 커널은 누구 생각이었을까요. 1992년 논문의 저자인 Isabelle Guyon이 2016년에 직접 쓴 회고에 따르면, 직선을 긋는 최대 마진 학습이 먼저 잘 되자, Vapnik은 특징끼리 곱해서 특징을 늘리자고 했고, Guyon은 대신 예전 potential function 방법의 커널 트릭을 쓰자고 했다고 합니다. 세 번째 시간에 배울 바로 그 커널입니다. Guyon은 Vapnik과 사무실을 같이 쓰며 열띤 토론을 했다고 적었고, 1저자 Bernhard Boser는 Guyon의 남편입니다. 다만 이것은 본인 회고라, 다른 사람의 기억과 다를 수 있습니다. 둘째, 1995년 같은 연구소에서 유명한 내기가 있었습니다. 그때 신경망은 숫자 인식도 잘했지만 왜 잘 되는지 설명하는 이론은 없었습니다. 연구소 그룹장 Larry Jackel은 5년 안에, 2000년까지는 그 이유를 알게 될 것이라고 걸었습니다. 증명되는 이론을 중시한 Vapnik은 이유도 모르는 큰 신경망은 2005년이면 아무도 역전파로 학습하지 않을 것이라고 걸었습니다. Yann LeCun도 같은 연구소에 있었고, Geoffrey Hinton이 강의에서 이 내기 문서를 소개한 적이 있습니다. 결과는 둘 다 졌습니다. Jackel이 건 것처럼 2000년까지 이유를 알아내지 못했고, Vapnik이 건 것처럼 신경망이 사라지지도 않았습니다. 신경망이 왜 그렇게 잘 되는지는 아직도 다 풀리지 않았고, 신경망은 지금 가장 많이 쓰는 모델이 됐습니다. 그래도 SVM은 2000년대 내내 가장 많이 쓰인 분류기 중 하나였고, 답이 하나로 정해지고 이론이 깔끔해서 지금도 데이터가 아주 크지 않을 때 많이 씁니다. 10년 뒤 2014년에는 Vapnik이 LeCun이 이끄는 페이스북 AI 연구소에 합류했습니다.}
\end{frame}

\section{초평면}
% --- 초평면
\begin{frame}{초평면 (hyperplane): 공간을 반으로 가르는 평평한 면}
\vfill
\centering\includegraphics[width=0.78\linewidth]{sv_hyper.png}\par\vskip0.3em
{\large 2차원이면 \kb{직선}, 3차원이면 \kb{평면}}
\vfill
\note{SVM은 초평면이라는 말에서 출발합니다. n차원 공간에서 초평면은 한 차원 낮은, 평평한 면입니다. 특징이 두 개인 2차원 평면에서는 직선이고, 특징이 세 개인 3차원 공간에서는 평면입니다. 특징이 더 많아져도 이름은 그대로 초평면입니다. 오늘은 그림을 그릴 수 있게 대부분 2차원, 즉 직선으로 보겠습니다.}
\end{frame}

% --- 초평면 식
\begin{frame}{초평면의 식: 값이 0 인 곳, 부호로 양쪽을 나눈다}
\vfill
\begin{columns}[c]
\begin{column}{0.44\textwidth}\centering\includegraphics[width=0.92\linewidth]{sv_half.png}\end{column}
\begin{column}{0.54\textwidth}
{\Large $w_1 x_1 + w_2 x_2 + b = 0$}\\[0.8em]
{\large 값 \textcolor{ksared}{$> 0$} 한쪽, \textcolor{ksablue}{$< 0$} 다른 쪽}\\[0.8em]
\begin{itemize}\large
\item 예: $1 + 2x_1 + 3x_2 = 0$
\begin{itemize}\normalsize
\item $(0.8, 0.6) \to \textcolor{ksared}{+4.4}$
\item $(-0.9, -0.7) \to \textcolor{ksablue}{-2.9}$
\end{itemize}
\end{itemize}
\end{column}
\end{columns}
\vfill
\note{2차원 초평면, 직선은 w1 x1 더하기 w2 x2 더하기 b가 0인 점들의 모임입니다. 예를 들어 1 더하기 2 x1 더하기 3 x2가 0인 직선을 그리면 왼쪽 검은 선입니다. 이 직선 위에 있지 않은 점을 넣으면 값이 0이 아니게 나오는데, 한쪽은 모두 양수, 다른 쪽은 모두 음수입니다. 점 0.8, 0.6을 넣으면 플러스 4.4, 점 마이너스 0.9, 마이너스 0.7은 마이너스 2.9입니다. 경계 가까이 있는 점은 0.2처럼 0에 가깝습니다. 그래서 초평면 하나가 공간을 반으로 나누고, 값의 부호만 보면 어느 쪽인지 압니다. 특징이 n개면 항이 n개로 늘어날 뿐 같습니다.}
\end{frame}

% --- 분리 초평면 분류기
\begin{frame}{분리 초평면으로 분류하기}
\vfill
\begin{columns}[c]
\begin{column}{0.36\textwidth}\centering\includegraphics[width=\linewidth]{sv_hook_d.png}\end{column}
\begin{column}{0.63\textwidth}
{\normalsize 학습 데이터 \; $(\mathbf{x}_1, y_1),\ \dots,\ (\mathbf{x}_p, y_p)$, \;$p = 20$}\\[0.2em]
{\normalsize $y_k = \begin{cases} +1 & \mathbf{x}_k \in \text{클래스 A (빨강)}\\ -1 & \mathbf{x}_k \in \text{클래스 B (파랑)} \end{cases}$}\\[0.5em]
{\large $D(\mathbf{x}) = \mathbf{w}\cdot\mathbf{x} + b = \sum_{i=1}^{n} w_i x_i + b$}\\[0.2em]
{\normalsize 분류: $D(\mathbf{x}) > 0$ 이면 A, 아니면 B}\\[0.5em]
{\normalsize 모든 점이 제 색 영역에 \;$\Leftrightarrow$\; \kb{$y_k\, D(\mathbf{x}_k) > 0$}}\\[0.2em]
{\normalsize 그런 선은 검은 선 · 회색 선처럼 \kq{무수히} 많다}
\end{column}
\end{columns}
\vfill
\note{이제 문제를 원 논문처럼 적어 봅시다. 학습 데이터는 점과 라벨의 짝 p개, 여기서는 20개입니다. 점 x k가 클래스 A, 빨강이면 라벨 y k는 플러스 1, 클래스 B, 파랑이면 마이너스 1입니다. 그리고 초평면의 식 w 점곱 x 더하기 b를 D라고 씁니다. SVM 원 논문, Boser, Guyon, Vapnik의 1992년 논문이 쓰는 이름이고, 결정 함수라는 뜻의 decision function에서 왔습니다. 빨간 점에서 D가 양수, 파란 점에서 음수면 잘 가른 것입니다. 두 조건을 한 줄로 쓰면, 라벨 곱하기 D가 모든 점에서 0보다 크다입니다. 라벨이 마이너스 1이면 음수 곱하기 음수라 양수가 되니까요. 새 점이 오면 D가 0보다 크면 A, 아니면 B로 분류하고, D의 절댓값이 클수록, 즉 경계에서 멀수록 더 확신합니다. 로지스틱회귀의 z가 0인 경계도 이런 분리 초평면입니다. 문제는 데이터가 깔끔하게 나뉘면 가르는 선이 그림처럼 무수히 많다는 것입니다. 어느 것을 골라야 할까요.}
\end{frame}

\section{최대 마진}
% --- 가장 넓은 길
\begin{frame}{최대 마진 분류기: 가장 넓은 길}
\markpage{street}
\vfill
\begin{columns}[c]
\begin{column}{0.44\textwidth}\centering\includegraphics[width=0.95\linewidth]{sv_street.png}\end{column}
\begin{column}{0.52\textwidth}
{\large \kb{마진 $M \in \mathbb{R}$} = 가장 가까운 점까지 거리}\\[0.8em]
{\large $M$ 이 가장 큰 선 = 길의 가운데}\\[0.8em]
{\large $M =$ \kb{0.67} \quad 동그라미 = \textbf{서포트 벡터}}
\end{column}
\end{columns}
\vfill
\note{가장 자연스러운 선택은 양쪽 점에서 가장 멀리 떨어진 초평면입니다. 경계에서 가장 가까운 학습 점까지의 거리를 마진, 대문자 M이라고 부르고, 마진이 가장 큰 초평면을 최대 마진 초평면, 이것으로 분류하는 것을 최대 마진 분류기라고 합니다. 그림으로는 두 무리 사이에 점이 하나도 없는 가장 넓은 길을 내고 그 가운데 선을 경계로 쓰는 것입니다. 이 데이터에서 M은 0.67, 길 전체 폭은 그 두 배인 1.34입니다. 경계까지의 거리가 정확히 M인 점, 점선에 닿은 노란 동그라미 점들을 서포트 벡터라고 부릅니다.}
\end{frame}

% --- 서포트 벡터만
\begin{frame}{왜 ``서포트 벡터'' 인가: 이 점들만 경계를 받친다}
\vfill
\centering\includegraphics[width=0.62\linewidth]{sv_move.png}\par\vskip0.2em
{\large 먼 점은 옮겨도 경계가 \kq{그대로}}
\vfill
\note{이 경계는 20개 점 중 단 3개, 서포트 벡터가 정합니다. 빨강 하나, 파랑 둘입니다. 길 가장자리에서 멀리 떨어진 점은 경계에 아무 영향이 없습니다. 오른쪽 그림에서 왼쪽 아래의 파란 점 하나를 더 멀리 옮겼는데, 경계와 길은 그대로입니다. 반대로 서포트 벡터 하나를 옮기면 경계가 바뀝니다. 로지스틱회귀는 모든 점이 조금씩 경계에 영향을 주는 것과 대조적입니다.}
\end{frame}


% --- 왜 2~3개로 정해지나
\begin{frame}{왜 점 2$\sim$3개로 정해지나: 닿을 때까지 넓힌다}
\vfill
\centering\includegraphics[width=0.94\linewidth]{sv_lock.png}\par\vskip0.2em
{\large 정할 것 3개 (기울기 · 위치 · $M$) $=$ 가장자리에 닿는 점 \kb{3개}}
\vfill
\note{왜 경계가 점 두세 개로 정해지는지 봅시다. 길을 최대로 넓혔다면 어떤 모양이어야 할까요. 첫째 그림처럼 한쪽, 빨강만 가장자리에 닿아 있으면 선을 파랑 쪽으로 평행이동하면 길이 넓어집니다. 그러니 최대가 아닙니다. 양쪽에 적어도 하나씩은 닿아야 합니다. 둘째 그림처럼 양쪽에 하나씩, 두 점만 닿아 있으면 선을 살짝 돌려서 두 점 모두에서 멀어질 수 있습니다. 예외는 선이 두 점을 잇는 선분에 수직일 때뿐이고, 그때는 두 점으로 끝납니다. 그 밖에는 돌리다 보면 세 번째 점이 닿아서 더 못 돌립니다. 셋째 그림이 그 상태, 우리 답입니다. 숫자로 보면 2차원에서 정할 것은 기울기, 위치, 그리고 M 세 개이고, 가장자리에 닿은 점 하나가 거리 = M이라는 식 하나를 주니 점 세 개면 딱 정해집니다. 특징이 n개면 보통 n 더하기 1개까지 닿습니다. 나머지 점은 M보다 멀어서 식에 들어가지 않고, 조금 움직여도 답이 바뀌지 않습니다.}
\end{frame}

% --- 2개 vs 3개
\begin{frame}{2개로 정해질 때 vs 3개로 정해질 때}
\vfill
\centering\includegraphics[width=0.94\linewidth]{sv_two_three.png}\par\vskip0.2em
{\large 가장 가까운 쌍의 수직이등분선 길이 \kb{비었으면 2개}, 점이 \kq{걸리면 3개}}
\vfill
\note{두 점으로 정해질 때와 세 점으로 정해질 때를 비교합니다. 왼쪽은 작은 예입니다. 빨강과 파랑 사이에서 가장 가까운 두 점을 잇고, 그 선분의 수직이등분선을 경계로 하면 M은 선분 길이의 절반, 0.71입니다. 이 길 안에 다른 점이 하나도 없으면 이것이 답이고, 두 점이 정합니다. 선이 두 점을 잇는 선분에 정확히 수직이라 어느 쪽으로 돌려도 한 점이 가까워지기 때문입니다. 가운데는 우리 데이터에 같은 방법을 쓴 것입니다. 가장 가까운 쌍은 빨강 0.7, 0.2와 파랑 마이너스 0.6, 마이너스 0.4이고 거리가 1.43이라 M은 0.72가 될 것 같지만, 그 길 안에 빨간 동그라미 점들이 걸립니다. 그래서 이 선은 쓸 수 없습니다. 오른쪽이 실제 답입니다. 파랑 두 점, 마이너스 0.6, 마이너스 0.4와 0.1, 마이너스 1.1을 이은 직선이 파랑 쪽 가장자리가 되어 기울기를 정합니다. 기울기가 정확히 마이너스 1이라 w가 0.71, 0.71입니다. 빨강 0.7, 0.2에서 그 직선까지의 거리 1.34가 길 폭 2M이 되어 폭을 정합니다. 그래서 M은 0.67로, 두 점으로 정할 때의 0.72보다 작습니다. 점이 많을수록 길 안에 점이 걸리기 쉬워서 세 점, 특징이 n개면 n 더하기 1개인 경우가 훨씬 흔합니다.}
\end{frame}

\begin{frame}{참고: 점에서 경계까지 거리 공식이 성립하는 이유}
\vfill
\begin{columns}[c]
\begin{column}{0.40\textwidth}\centering\includegraphics[width=\linewidth]{sv_cos.png}\end{column}
\begin{column}{0.59\textwidth}
{\large \textcolor{fdist}{\textbf{점 $\mathbf{x}$ 에서 경계선까지 거리} $= \dfrac{|D(\mathbf{x})|}{\lVert\mathbf{w}\rVert}$}}\\[0.6em]
\hspace*{2.5em}\fcolorbox{fbeta}{fbeta!8}{\small ※ 두 벡터 사이 각: $\cos\theta = \dfrac{\mathbf{a}\cdot\mathbf{b}}{\lVert\mathbf{a}\rVert\,\lVert\mathbf{b}\rVert}$}\\[0.5em]
{\normalsize \color{fsh}\textbf{1.} $\mathbf{x}$ 의 위치 $\lVert\mathbf{x}\rVert\cos\theta = \tfrac{\mathbf{w}\cdot\mathbf{x}}{\lVert\mathbf{w}\rVert}$}\\[0.4em]
{\normalsize \color{fbd}\textbf{2.} 경계선 위의 점 $\mathbf{z}$ 는 $D(\mathbf{z}) = 0$, 곧 $\mathbf{w}\cdot\mathbf{z} = -b$\\ \hspace*{1.2em}$\to$ 경계의 위치 $= \tfrac{-b}{\lVert\mathbf{w}\rVert}$}\\[0.4em]
{\normalsize \textbf{3.} \textcolor{fdist}{거리} $= \tfrac{\mathbf{w}\cdot\mathbf{x}}{\lVert\mathbf{w}\rVert} - \tfrac{-b}{\lVert\mathbf{w}\rVert} = \tfrac{\mathbf{w}\cdot\mathbf{x} + b}{\lVert\mathbf{w}\rVert} = \tfrac{D(\mathbf{x})}{\lVert\mathbf{w}\rVert}$}\\[0.5em]
{\normalsize 예 $\mathbf{w} = (1, 1)$, $b = 0.5$: \;\textcolor{fsh}{1.41} $-$ (\textcolor{fbd}{$-0.35$}) $=$ \textbf{\textcolor{fdist}{1.77}}}
\end{column}
\end{columns}
\vfill
\note{마진 M은 점에서 경계선까지의 거리입니다. 점과 점 사이가 아니라, 점에서 경계선에 수선을 내린 길이, 그림의 초록 화살표입니다. 이 거리는 D x의 절댓값을 w의 길이로 나눈 것인데, 왜 그런지 코사인으로 세 줄이면 됩니다. 첫째, 두 벡터 사이의 코사인은 점곱을 두 길이로 나눈 것입니다. 양변에 x의 길이를 곱하면, x의 길이 곱하기 코사인 세타, 즉 점 x에서 w 축으로 수선을 내린 발까지의 길이, x의 w 방향 위치가 w 점곱 x 나누기 w의 길이입니다. 그림의 파란 치수선, 1.41입니다. 둘째, 경계는 w에 수직인 직선이라 경계 위의 점은 모두 위치가 같습니다. 경계선 위의 점을 z라고 하면 D z가 0이니 w 점곱 z가 마이너스 b이고, 그래서 경계의 위치는 마이너스 b 나누기 w의 길이입니다. 그림의 회색, 마이너스 0.35입니다. 셋째, 점에서 경계까지의 거리는 두 위치의 차이입니다. 빼면 w 점곱 x 더하기 b를 w의 길이로 나눈 것, 즉 D x 나누기 w의 길이가 됩니다. 반대편 점은 음수가 나오니 절댓값을 씌웁니다. w의 길이가 1일 필요는 없습니다. 나누기 w의 길이는 코사인 식에서 원래 나오는 것입니다. 예는 w가 1, 1이고 b가 0.5, 경계 x1 더하기 x2가 마이너스 0.5인 직선이고, 점은 0.4, 1.6입니다. w의 길이는 루트2, 1.41입니다. 점의 위치는 w 점곱 x 2 나누기 1.41, 1.41이고, 경계의 위치는 마이너스 0.5 나누기 1.41, 마이너스 0.35입니다. 차이는 1.77입니다. D x는 0.4 더하기 1.6 더하기 0.5, 2.5이고, 2.5 나누기 1.41도 1.77로 같습니다. 같은 선을 2 x1 더하기 2 x2 더하기 1로 쓰면 D는 5가 되지만 w의 길이도 2.83이 되어 거리는 똑같이 1.77입니다. D는 식을 몇 배로 쓰느냐에 따라 달라지고, w의 길이로 나누면 그 배율이 사라집니다. 앞으로 학습할 때는 w의 길이를 1로 맞춰서, 나누기 없이 D가 곧 거리가 되게 씁니다.}
\end{frame}

% --- 참고: 점곱과 코사인 (값 = 그림자 - 경계의 그림자 = 거리)
\begin{frame}{값에 $y_k$ 를 곱한다: 파랑 $D(\mathbf{x}) \to -D(\mathbf{x})$}
\vfill
\centering
\includegraphics[width=0.68\textwidth]{sv_fold.png}\\[0.1em]
{\large 빨강 $D(\mathbf{x}_k)$, \; 파랑 $-D(\mathbf{x}_k)$ \;$\to$\; 둘 다 $y_k D(\mathbf{x}_k)$ = 거울까지 거리}\\[0.3em]
{\large 모든 점이 경계에서 $M$ 이상: \;$\dfrac{y_k D(\mathbf{x}_k)}{\lVert\mathbf{w}\rVert} \ge M$}
\vfill
\note{이번에는 점이 아니라 값, w 점곱 x 더하기 b에 y k를 곱합니다. 앞에서 점에 y k를 곱한 것, 원점 기준 180도 회전과는 다른 그림입니다. 파란 점은 y가 마이너스 1이라 값의 부호가 뒤집힙니다. 값은 경계까지의 거리이니, 부호만 뒤집힌 값을 갖는 자리는 경계 반대편으로 같은 거리, 즉 경계를 거울 삼아 비춘 자리입니다. 그래서 값에 y를 곱하는 것은 파란 점을 경계에 비춰 빨강 쪽으로 넘긴 것과 같습니다. 원점 기준으로 180도 돌리면 b가 0이 아닐 때 값이 달라집니다. 접고 나면 20개가 모두 한쪽에 있고, 각 점의 값은 거울까지의 거리입니다. 틀린 쪽에 있던 점이 있다면 접어도 거울 반대편에 남아 값이 음수가 됩니다. 빨강은 값 그대로, 파랑은 값에 마이너스를 붙이면 둘 다 거울까지의 거리가 되는데, 이 둘을 한 식으로 쓴 것이 y k 곱하기 D x k입니다. 그중 가장 작은 것이 M이고, 거울을 기울이고 옮겨서 이 M을 최대로 만드는 것이 바로 SVM, 최대 마진입니다. 뒤의 식 쪽에서 이 문장이 그대로 조건 y k D x k ≥ M이 됩니다. 답에서는 빨강 한 점과 비춰 넘긴 파랑 두 점이 똑같이 0.67 거리에 닿습니다. 거울을 움직이면 접힌 모양도 같이 바뀌니, 가장 가까운 점도 계속 바뀝니다.}
\end{frame}

\section{식}
% --- 최대 마진 식 해설 (말풍선)
\begin{frame}{최대 마진 초평면 찾기: 한 줄씩 읽기}
\vfill
\begin{minipage}{0.62\textwidth}
{\Large\begin{align*}
\text{find} \quad & w_1, w_2,\, b\\ & \text{that maximize } \tikzmarknode{mm}{M > 0}\\[0.4em]
\text{subject to} \quad & \tikzmarknode{nb}{w_1^2 + w_2^2 = 1},\\[0.4em]
& \tikzmarknode{cc}{y_k(w_1 x_{k1} + w_2 x_{k2} + b) \ge M}
\end{align*}}
\end{minipage}
\begin{tikzpicture}[remember picture,overlay]
\node[bub,anchor=west] (b1) at ([xshift=1.4cm,yshift=0.55cm]mm.east) {마진 $M$ 최대};
\draw[bubarr] (b1.west) -- (mm.east);
\node[bub,anchor=west] (b2) at ([xshift=1.4cm,yshift=-0.1cm]nb.east) {$\mathbf{w}$ 의 길이 $\lVert\mathbf{w}\rVert = \sqrt{w_1^2 + w_2^2}$ 를 1 로};
\draw[bubarr] (b2.west) -- (nb.east);
\node[bub,anchor=west] (b3) at ([xshift=0.6cm,yshift=-0.4cm]cc.east) {모든 점이 $M$ 이상};
\draw[bubarr] (b3.west) -- (cc.east);
\end{tikzpicture}
\vskip1.2em
\centering
{\large 이 데이터: \; $\mathbf{w} = (0.71,\ 0.71)$, \; $b = 0.04$, \; \kb{$M = 0.67$}}
\vfill
\note{한 줄씩 읽어 봅시다. 첫 줄, 마진 M을 최대로 합니다. 마진은 경계에서 가장 가까운 점까지의 거리입니다. 둘째 줄, w1 제곱 더하기 w2 제곱이 1, 즉 w의 L2 norm 제곱이 1, 길이 L2 norm을 1로 맞춥니다. b은 들어가지 않습니다. 그래야 식의 값이 곧 경계까지의 거리가 됩니다. 이 조건이 없으면 w를 두 배, 세 배 키우기만 해도 값이 커지니 M을 끝없이 키울 수 있어서 문제가 성립하지 않습니다. 점과 직선 사이 거리 공식에서 거리는 값의 절댓값을 w의 길이로 나눈 것이라, 길이를 1로 맞추면 값이 그대로 거리입니다. 이 공식이 왜 성립하는지는 앞의 참고 쪽에서 봤습니다. 셋째 줄, 모든 점에서 라벨 곱하기 값이 M 이상, 즉 모든 점이 자기 쪽으로 경계에서 M 이상 떨어져 있어야 합니다. 이 데이터의 답은 w가 0.71, 0.71, b가 0.04, M이 0.67입니다. 서포트 벡터 0.7, 0.2을 넣으면 값이 0.67, 딱 M입니다. 이 문제는 볼록 최적화라서 효율적으로 풀 수 있고 답이 하나로 정해집니다.}
\end{frame}

% --- 최대 마진 식 (공식 그대로, ISL 9.9-9.11)
\begin{frame}{최대 마진 초평면 찾기: Generalize ($n$ 개 특징, $p$ 개 점)}
\markpage{form}
\vfill
\centering
{\large\begin{align*}
\text{find} \quad & w_1, \dots, w_n,\, b\\ & \text{that maximize } M > 0\\
\text{subject to} \quad & \sum_{i=1}^{n} w_i^2 = 1,\\
& y_k(w_1 x_{k1} + w_2 x_{k2} + \cdots + w_n x_{kn} + b) \ge M \quad \text{for all } k = 1, \dots, p
\end{align*}}
{\large $y_k = +1$ (A) 또는 $-1$ (B)}

\vfill
\note{최대 마진 초평면을 찾는 문제를 원 논문, Boser, Guyon, Vapnik 1992년 논문의 식 9 그대로 쓰면 이렇습니다. 학습 데이터가 p개, 특징이 n개이고 라벨은 클래스 A면 플러스 1, B면 마이너스 1입니다. 점 번호는 k, 특징 번호는 i로 씁니다. w1부터 wn과 b, 그리고 마진 M을 골라서 M을 최대로 합니다. 조건은 두 가지입니다. w1부터 wn까지의 제곱합이 1, 그리고 모든 점에서 라벨 곱하기 초평면 식의 값이 M 이상. 앞 쪽에서 읽은 2차원 식을 특징이 n개인 경우로 그대로 늘린 것입니다.}
\end{frame}

\section{학습}
% --- 학습 0: 숫자 3개 찾기
\begin{frame}{푸는 식: 가장 가까운 점을 경계에서 가장 멀리}
\vfill
\centering
\vspace*{2.1cm}
{\LARGE $M^* = \ \tikzmarknode{mx}{\displaystyle\max_{\mathbf{w},\, b,\ \lVert\mathbf{w}\rVert = 1}}\ \ \ \tikzmarknode{mn}{\displaystyle\min_{k}}\ \ \ \tikzmarknode{yd}{y_k D(\mathbf{x}_k)}$}
\begin{tikzpicture}[remember picture,overlay]
\node[bub,anchor=south west] (b1) at ([xshift=-1.15cm,yshift=0.6cm]yd.north east) {\textbf{점 $\mathbf{x}_k$ 에서 경계까지 거리}\\ ($\lVert\mathbf{w}\rVert = 1$ 이라 값 = 거리)};
\draw[bubarr] ([xshift=-0.6cm]b1.south) -- (yd.north);
\node[bub,anchor=south] (b2) at ([xshift=-1.2cm,yshift=0.7cm]mn.north) {\textbf{점 $p$ 개의 거리 중 가장 작은 것}\\ = 가장 가까운 점의 거리 = 마진 $M$};
\draw[bubarr] (b2.south) -- (mn.north);
\node[bub,anchor=north] (b3) at ([xshift=0.4cm,yshift=-0.6cm]mx.south) {\textbf{선을 돌리고 ($\mathbf{w}$) 옮기며 ($b$) 그 거리를 가장 크게}\\ $\lVert\mathbf{w}\rVert = 1$: $\mathbf{w}$ 길이 고정 (안 묶으면 값이 끝없이 커진다)};
\draw[bubarr] (b3.north) -- (mx.south);
\end{tikzpicture}
\vskip2.3cm
{\normalsize 예: 평균 선은 가장 가까운 점까지 0.62, \;SVM 선은 \kb{0.67} $\to$ 더 큰 SVM 선이 답}\\[0.3em]
{\normalsize 방법: \kq{가장 가까운 점}이 민다 (min) $\to$ 그 거리가 커진다 (max) $\to$ $\mathbf{w}$ 길이 1 로 $\to$ 반복}
\vfill
\note{방금 본 M 최대, 모든 점이 M 이상이라는 문제를 원 논문은 한 줄로도 씁니다. 선 하나를 정하면 모든 점이 M 이상이어야 하니 M은 가장 가까운 점의 거리까지만 커질 수 있습니다. 그래서 M 자리에 가장 작은 거리, min을 바로 쓸 수 있습니다. M 별은 모든 선 중에서 가장 큰 마진이고, 이 데이터에서는 0.67입니다. 거리가 M 별인 점이 서포트 벡터입니다. max min y k D x k입니다. 안쪽부터 읽습니다. 첫째, y k D x k는 점 k의 값에 라벨을 곱한 것이고, w의 길이를 1로 묶어 두었으니 점 k에서 경계까지의 거리입니다. 틀린 쪽에 있는 점이면 음수가 됩니다. 둘째, min k는 점 20개의 거리 중 가장 작은 것, 즉 경계에 가장 가까운 점까지의 거리를 고릅니다. 이것이 마진 M입니다. 셋째, 바깥의 max는 w와 b를 바꿔 가며, 즉 경계를 돌리고 옮겨 가며 그 가장 작은 거리를 가장 크게 만듭니다. 그래서 전체는 가장 가까운 점을 경계에서 가장 멀리 두는 선을 찾으라는 뜻이고, 앞에서 본 가장 넓은 길과 같은 말입니다. 원 논문은 max 아래에 w만 적지만, b도 같이 바꿔 가며 고르기 때문에 여기서는 b도 적었습니다. 아래 첨자에 w의 길이가 1이라고 적은 것은, 이 조건이 없으면 w를 두 배, 세 배 키우기만 해도 값이 커져서 max가 끝없이 커지기 때문입니다. 우리가 정할 숫자는 w1, w2, b 셋이고, 방법은 다음 쪽부터 봅니다. 가장 가까운 점이 선을 밀고, w의 길이를 다시 1로 맞추는 것을 반복합니다. 가장 가까운 점이 미는 것이 min이고, 밀어서 그 거리를 키우는 것이 max입니다.}
\end{frame}

% --- 출발: 0걸음 → 1걸음
\begin{frame}{학습 (max min, $\lVert\mathbf{w}\rVert = 1$): $\mathbf{w} = 0$ 에서 첫 선}
\vfill
\begin{columns}[c]
\begin{column}{0.42\textwidth}\centering\includegraphics[width=0.95\linewidth]{sv_step1.png}\end{column}
\begin{column}{0.56\textwidth}
\renewcommand{\arraystretch}{1.6}
{\normalsize\begin{tabular}{l}
\textbf{0걸음}: $\mathbf{w} = (0, 0)$, $b = 0$\\
$\to$ 점 20개의 값이 \kq{모두 0} $\to$ 20개가 다 같이 민다\\[0.3em]
\textbf{1걸음 뒤}: $\mathbf{w} \propto \sum_{k=1}^{p} y_k\mathbf{x}_k = (16.5,\ 18.2)$\\
$\to$ 길이 1 로: $\mathbf{w} = (0.67,\ 0.74)$, \;$b = 0$\\[0.3em]
= 앞에서 본 \kb{평균 선}과 같은 선
\end{tabular}}
\end{column}
\end{columns}
\vfill
\note{이제 학습입니다. 방금 본 max min 식, w의 길이가 1인 식을 풉니다. 이 방법은 이해하기 쉽게 고른 수업용 방법이고, 원 논문이 쓴 방법, 알파로 푸는 것은 다음 시간에 봅니다. 가장 가까운 점이 미는 것이 min이고, 밀어서 그 거리를 키우는 것이 max입니다. 출발점을 따로 봅시다. 처음에는 w가 0, 0이고 b가 0입니다. 이때는 어느 점을 넣어도 값이 0이라, 20개가 모두 똑같이 가장 가까운 점이 되고, 첫 걸음에는 20개가 모두 밉니다. 그래서 w는 학습률 곱하기 y k x k의 합, 빨강 10개 좌표 합 8.3, 9.1에서 파랑 합 마이너스 8.2, 마이너스 9.1을 뺀 16.5, 18.2 방향이 됩니다. b에는 학습률 곱하기 y의 합을 더하는데, 빨강 10개와 파랑 10개라 y의 합이 0이고 b는 0 그대로입니다. 길이를 1로 맞추면 w는 0.67, 0.74입니다. 이 방향은 20개 y k x k의 평균 방향과 같으니, 첫 선은 앞에서 본 가장 쉬운 방법, 평균 선과 똑같습니다. 그림이 그 첫 선이고, b가 0이라 원점을 지납니다. 다음 쪽부터는 이 선에서 가장 가까운 점이 한 걸음씩 밉니다.}
\end{frame}

% --- 학습 한 걸음
\begin{frame}{학습 한 걸음: 가장 가까운 점이 민다}
\markpage{train}
\vfill
\centering
\renewcommand{\arraystretch}{1.6}
{\large\begin{tabular}{rl}
\textbf{1.} & 점마다 거리 $y_k D(\mathbf{x}_k)\,/\,\lVert\mathbf{w}\rVert = y_k D(\mathbf{x}_k)$ \;($\lVert\mathbf{w}\rVert = 1$ 이라)\\
\textbf{2.} & 값이 \kq{가장 작은 점} $k = \arg\min_l\, y_l D(\mathbf{x}_l)$ 가 민다:\\
 & $\mathbf{w}' = \mathbf{w} + \eta\, y_k \mathbf{x}_k$, \quad $b' = b + \eta\, y_k$\\
\textbf{3.} & $\mathbf{w} \leftarrow \mathbf{w}'/\lVert\mathbf{w}'\rVert$ \;(길이 1 로), \quad $b \leftarrow b'$ \;(그대로)
\end{tabular}}\\[0.8em]
{\large 밀고 나면 그 점의 값: \; $y_k(\mathbf{w}'\cdot\mathbf{x}_k + b') = y_k(\mathbf{w}\cdot\mathbf{x}_k + b) + \eta\,(\lVert\mathbf{x}_k\rVert^2 + 1)$}\\[0.6em]
{\normalsize 학습률 $\eta = 0.05/\sqrt{\text{걸음}}$}
\vfill
\note{출발은 w가 0, 0, b가 0입니다. 그러면 어느 점을 넣어도 값이 0이라 20개가 모두 똑같이 가장 작은 점이 되고, 첫 걸음에는 20개가 모두 밉니다. 그래서 첫 w는 학습률 곱하기 20개의 y x 합, b는 학습률 곱하기 y의 합, 0입니다. 이것은 20개 y x의 평균에 학습률 곱하기 20을 곱한 것이라 방향이 똑같고, b도 0이니 첫 선은 앞에서 본 가장 쉬운 방법, 모든 벡터의 평균 선과 똑같습니다. 숫자로는 빨강 10개 좌표 합 8.3, 9.1에서 파랑 합 마이너스 8.2, 마이너스 9.1을 빼면 16.5, 18.2, 학습률 0.05를 곱하면 0.83, 0.91, 길이 1.23으로 나누면 0.67, 0.74입니다. 그다음부터 한 걸음은 세 단계입니다. 첫째, 점 20개를 넣어 값, 라벨 곱하기 w x 더하기 b를 계산합니다. w의 길이가 1이니 이 값이 곧 경계까지의 거리이고, 틀린 점은 음수입니다. 둘째, 값이 가장 작은 점 하나가 밉니다. 빨간 점이면 w에 자기 좌표를 학습률만큼 더하고 b를 올리고, 파란 점이면 좌표를 빼고 b를 내립니다. 그러면 그 점의 값이 커집니다. 실제로 바꾼 w, b를 넣어 보면 y k 제곱이 1이라서 값이 학습률 곱하기 (좌표 길이 제곱 더하기 1)만큼 늘어납니다. 빨강이든 파랑이든 항상 양수입니다. 즉 가장 가까운 점이 경계를 자기 반대쪽으로 밀어내는 것입니다. 우리가 키우려는 M이 바로 이 가장 작은 값이니, 그 점을 키우는 것이 M을 키우는 방향입니다. 셋째, w만 길이로 나눠 다시 1로 맞추고, b는 그대로 둡니다. 길이 1 조건은 w에만 걸려 있기 때문입니다. b까지 나누면 매 걸음 b가 0 쪽으로 조금씩 끌려가서, 답의 b가 큰 데이터에서는 답에 못 갑니다. 무작위 데이터 200개로 sklearn과 비교해 보면 b까지 나눌 때는 17개가 답에서 벗어나고, w만 나누면 200개 모두 sklearn과 같은 선에 닿습니다. 점들을 z-score 해 두었으니 w와 b에 같은 학습률 하나를 씁니다. 학습률은 걸음이 갈수록 줄입니다. 최솟값을 키우는 방향으로 조금씩 가는 것이라 경사하강법의 한 종류, 정확히는 서브그래디언트 방법입니다.}
\end{frame}

\begin{frame}{한 번 밀면: $\mathbf{w}$ 는 회전, $b$ 는 평행이동}
\vfill
\centering
\includegraphics[width=0.82\textwidth]{sv_push_dir.png}\\[0.2em]
{\large ①② 는 \kq{동시에} (나눠 그린 것) \quad 점 $\mathbf{x}_k$ 까지 거리: 0.618 $\to$ $\mathbf{w}$ 만 0.623 $\to$ 둘 다 0.659}
\vfill
\note{1걸음 뒤, 빨간 점 0.7, 0.2이 가장 가까울 때 한 번 미는 것을 식 순서대로 둘로 나눠 그렸습니다. 실제로는 w와 b를 한 걸음에 동시에 바꿉니다. w 업데이트에는 b가, b 업데이트에는 w가 들어가지 않으니 순서는 상관없고, 이해하기 쉽게 나눠 그린 것입니다. 실제 변화가 작아서 화살표는 10배로 그렸습니다. 왼쪽, w에 학습률 곱하기 y k 곱하기 x k를 더합니다. x k는 원점에서 점까지의 화살표이고, 빨간 점이라 y k가 플러스 1이니 x k를 더합니다. 회색 w는 선에 수직인 법선인데, 빨강 쪽을 가리킵니다. 여기에 x k를 더하면 검은 w 프라임이 x k 쪽으로 돌아가고, 법선이 돌았으니 선도 돕니다. 왜 거리가 늘어날까요. 앞의 참고 쪽에서 w 점곱 x는 x의 길이 곱하기 코사인 세타, 즉 x를 w 방향으로 비춘 길이였습니다. 빨간 점이면 w가 x k 쪽으로 돌아 이 길이가 길어지고, 파란 점이면 w가 x k에서 멀어지며 이 길이가 더 음수가 됩니다. 어느 쪽이든 점 k가 선에서 멀어집니다. 이번에는 거리가 0.618에서 0.623으로 조금 늡니다. 오른쪽, b에 학습률 곱하기 y k를 더하면 기울기는 그대로이고 그 선이 자기에 수직인 방향으로, 점 k 반대쪽으로 평행이동합니다. 거리는 0.659가 됩니다. 이처럼 점에서 직각으로만 미는 것이 아니라, 그 점의 값을 가장 빨리 키우는 방향으로 w와 b를 함께 바꾸는 것입니다.}
\end{frame}

% --- 걸음별
\begin{frame}{반복하면 미는 점이 3개만 남는다 (1/3)}
\vfill
\centering\includegraphics[width=0.95\linewidth]{sv_train9_1.png}\par\vskip0.2em
{\large 첫 선 = 평균 방향 $\to$ 가장 가까운 점이 밀 때마다 식이 조금씩 바뀐다}
\vfill
\note{1걸음 뒤 첫 선은 0.67 x1 더하기 0.74 x2 = 0, 모든 벡터의 평균 방향이고 원점을 지납니다. M은 0.62이고 빨간 점 0.7, 0.2가 가장 가까워 다음에 밉니다. 빨간 점이 밀면 w에 그 좌표가 더해져 0.68, 0.73으로 조금 돌고 b가 0.04로 커집니다. 2걸음 뒤 M은 0.66. 같은 점이 한 번 더 밀어 3걸음 뒤에는 b가 0.06이 되는데, 이번에는 너무 밀어서 반대편 파란 점 마이너스 0.6, 마이너스 0.4가 가장 가까워집니다. M은 0.64로 오히려 조금 줄었습니다.}
\end{frame}

\begin{frame}{반복하면 미는 점이 3개만 남는다 (2/3)}
\vfill
\centering\includegraphics[width=0.95\linewidth]{sv_train9_2.png}\par\vskip0.2em
{\large 빨강 · 파랑이 \kq{번갈아} 민다, 10걸음 뒤 세 번째 서포트 벡터}
\vfill
\note{4걸음 뒤에는 파란 점이 밀어 b가 0.04로 돌아오고 M이 0.66, 파란 점이 한 번 더 밀어 5걸음 뒤에는 b가 0.02, M 0.65. 이렇게 빨강과 파랑의 가장 가까운 점이 번갈아 밀며 b가 0.02와 0.06 사이를 오갑니다. 10걸음 뒤에는 세 번째 서포트 벡터, 파란 점 0.1, 마이너스 1.1이 처음으로 가장 가까워집니다. w가 0.71, 0.70으로 거의 대각선 방향입니다.}
\end{frame}

\begin{frame}{반복하면 미는 점이 3개만 남는다 (3/3)}
\vfill
\centering\includegraphics[width=0.95\linewidth]{sv_train9_3.png}\par\vskip0.2em
{\large 식이 멈춘다 $\to$ 미는 점은 끝까지 \kb{서포트 벡터 3개}뿐}
\vfill
\note{100걸음 뒤에는 식이 0.71 x1 더하기 0.71 x2 더하기 0.04 = 0, 소수 둘째 자리까지는 이미 답과 같고 M도 0.67입니다. 1000걸음, 2만 걸음 뒤에도 식은 그대로이고, 서포트 벡터 3개가 번갈아 아주 조금씩 밀며 길을 붙잡고 있습니다. 1걸음 뒤로 한 번이라도 민 점은 이 3개뿐이고, 나머지 17개는 한 번도 밀지 않았습니다. 학습률이 걸음마다 줄어서 흔들림이 점점 작아집니다.}
\end{frame}

% --- 걸음마다 M
\begin{frame}{걸음마다 $M$ 이 어떻게 바뀌나}
\vfill
\centering\includegraphics[width=\linewidth]{sv_trainpath.png}
\vfill
\note{같은 학습을 걸음마다 M, 가장 가까운 점까지의 거리로 그렸습니다. 가로축은 로그 눈금이고, 세로축은 0.55에서 0.70까지만 크게 키워 그렸습니다. 1걸음 뒤 0.62에서 시작해, 가장 가까운 점이 밀 때마다 반대편 점이 조금 가까워지니 M이 위아래로 흔들립니다. 학습률이 걸음마다 줄면서 흔들림이 작아지고, 수백 걸음이면 SVM 답 0.67에서 거의 멈춥니다.}
\end{frame}

% --- 숫자로 따라가기
\begin{frame}{학습을 숫자로 따라가 보면}
\vfill
\centering
\renewcommand{\arraystretch}{1.25}
{\large\begin{tabular}{r|rrr|r|l}
걸음 뒤 & $w_1$ & $w_2$ & $b$ & $M$ & 다음에 미는 점\\ \hline
0 & 0 & 0 & 0 & 0 & 20개 모두 \; (선 없음, 값 모두 0)\\
1 & $0.67$ & $0.74$ & 0 & 0.62 & \textcolor{ksared}{$(0.7, 0.2)$}\\
2 & $0.68$ & $0.73$ & $0.04$ & 0.66 & \textcolor{ksared}{$(0.7, 0.2)$}\\
3 & $0.69$ & $0.72$ & $0.06$ & 0.64 & \textcolor{ksablue}{$(-0.6, -0.4)$}\\
10 & $0.71$ & $0.70$ & $0.04$ & 0.67 & \textcolor{ksablue}{$(0.1, -1.1)$}\\
2만 & $0.71$ & $0.71$ & $0.04$ & \kb{0.67} & 서포트 벡터 3개가 번갈아\\ \hline
SVM 답 & $0.71$ & $0.71$ & $0.04$ & 0.67 &
\end{tabular}}
\vfill
\note{실제로 점 20개로 돌린 숫자입니다. 각 줄은 그만큼 걸은 뒤의 숫자와, 그 다음 걸음에 미는 점입니다. 0걸음 뒤, 즉 출발할 때는 w, b가 모두 0이라 선이 없고 모든 점의 값이 0이라 M도 0입니다. 그래서 20개가 모두 밀어 첫 선이 생깁니다. 1걸음 뒤에는 빨간 점 0.7, 0.2이 가장 가까워 M이 0.62입니다. 이 점이 두 번 밀고, 3걸음 뒤에는 반대편 파란 점 마이너스 0.6, 마이너스 0.4이 가장 가까워집니다. 이렇게 가장 가까운 점이 번갈아 밀면서 M이 0.62와 0.67 사이에서 조금씩 흔들리다가, 10걸음 뒤에는 세 번째 서포트 벡터 파란 점 0.1, 마이너스 1.1까지 밀기 시작하고 거의 답에 닿습니다. 2만 걸음 뒤 w는 0.71, 0.71, b는 0.04, M은 0.67로 sklearn이 낸 SVM 답과 같습니다.}
\end{frame}

\begin{frame}{z-score 를 안 하면: 학습이 제대로 이루어지지 않는다}
\vfill
\centering
{\large 원래 값에 $x_2 + 100$ 만 했을 때}\\[0.2em]
\includegraphics[width=0.74\textwidth]{sv_zscore.png}\\[0.1em]
{\large $\mathbf{x}_k$ 가 모두 위쪽 화살표 + $b$ 는 $-82.8$ 까지 가야 함 $\to$ 학습이 제대로 안 된다}
\vfill
\note{왜 처음에 z-score를 했는지 봅시다. w에 x k를 더하는 학습이라 데이터가 원점에서 멀면 문제가 생깁니다. z-score 하기 전 원래 값의 x2에 100을 더해 봤습니다. 답은 그대로 100만큼 옮겨진 선이어야 합니다. 그런데 그대로 학습하면 2만 걸음 뒤에도 빨간 점 10개 중 9개가 틀립니다. 이유는 둘입니다. 첫째, 모든 점의 좌표가 (5, 105)처럼 거의 같은 위쪽 화살표라서 w를 어느 쪽으로 돌릴지 점마다 차이가 거의 없고 크게 휘둘립니다. 둘째, 답의 b가 마이너스 82.8인데 b는 한 걸음에 학습률만큼씩만 움직여 거기까지 못 갑니다. 오른쪽은 우리가 쓰는 z-score 한 점으로 같은 학습을 한 것으로, SVM 답 M 0.67에 도착합니다. 원래 값에 100을 더해도 z-score 하면 평균을 빼니 같은 점이 됩니다. 최대 마진 답 자체는 데이터를 옮겨도 같이 옮겨질 뿐인데, 이 학습 방법이 원점 위치와 크기에 민감한 것입니다. 단, 이 데이터처럼 축마다 표준편차가 다르면, 1과 3이면, 나누는 순간 거리 재는 기준이 바뀌어 답의 선도 달라집니다. z-score 해서 구한 답을 원래 값으로 되돌리면, 원래 값에서 바로 구한 최대 마진 선과 같지 않습니다. 어느 축이든 같은 크기로 보겠다는 선택을 한 것입니다.}
\end{frame}

% --- 학습 3: 가려내기
\begin{frame}{끝나고 거리를 재면 서포트 벡터가 보인다}
\vfill
\centering\includegraphics[width=0.78\linewidth]{sv_slack_bars.png}\par\vskip0.2em
{\large 거리 $= M$ 인 3개 = \kb{서포트 벡터}}
\vfill
\note{학습이 끝난 뒤 점 20개의 경계까지 거리를 작은 순서로 막대로 그렸습니다. 세 점이 모두 0.67로 M, 길 가장자리 위에 있습니다. 이 셋이 서포트 벡터입니다. 나머지 17개는 모두 0.81 이상입니다. 서포트 벡터를 미리 고르는 게 아니라, 가장 가까운 점이 미는 학습을 하다 보면 끝까지 가장 가까운 점들이 저절로 서포트 벡터가 됩니다. 실무에서 쓰는 sklearn의 SVC는 이런 방법 대신 SMO라는 다른 풀이법을 쓰지만 같은 답을 내고, support vectors라는 속성에 이 3개를 담아 줍니다.}
\end{frame}

% --- 가장 가까운 점 고르기 = 점마다 무게
\begin{frame}{``가장 가까운 점 고르기'' $=$ 점마다 \textbf{무게}를 나눠 주기}
\vfill
\centering
{\large 오늘의 식: \quad $\displaystyle \max_{\mathbf{w},\,b,\ \lVert\mathbf{w}\rVert = 1}\ \min_{k}\ y_k D(\mathbf{x}_k)$ \qquad 안쪽 $\min_k$ $=$ 가장 가까운 점 하나 고르기}\\[0.5em]
{\Large $\displaystyle \underbrace{\min_{k}\ y_k D(\mathbf{x}_k)}_{\text{\normalsize 최소값: 한 점}} \;=\; \underbrace{\min_{\lambda}\ \sum_k \lambda_k\, y_k D(\mathbf{x}_k)}_{\text{\normalsize 가중 평균의 최소: 여러 점이 나눠 가질 수 있다}}$}\\[0.1em]
{\small $\lambda_k \ge 0$, \;$\sum_k \lambda_k = 1$}\\[0.5em]
{\normalsize $\mathbf{w}$ 가 정해지면: $\lambda$ 는 \textbf{argmin 점 하나에 $1$}}\\[0.1em]
{\normalsize 답에서는 \kq{argmin 이 여럿} (서포트 벡터) $\to$ \textbf{나눠 가진다}}\\[0.4em]
\renewcommand{\arraystretch}{1.15}
{\large\begin{tabular}{c|ccc|c|c}
 & \textcolor{ksared}{빨강 $r$} & \textcolor{ksablue}{파랑 $b_1$} & \textcolor{ksablue}{파랑 $b_2$} & 나머지 17개 & 합\\ \hline
$\lambda$ & \kb{0.5} & \kb{0.25} & \kb{0.25} & 0 & 1
\end{tabular}}\\[0.4em]
{\large 이 무게가 학습한 $\mathbf{w}$ 안의 \textbf{몫}과 같다}
\vfill
\note{오늘의 식은 w의 길이를 1로 묶은 max min입니다. 안쪽 min은 점들 중 값이 가장 작은 점, 가장 가까운 점 하나를 고르는 것입니다. 이 일은 점마다 무게를 나눠 주고 무게 준 합을 가장 작게 만드는 일과 같습니다. 무게는 0 이상이고 다 더하면 1입니다. 무게 준 합은 값들의 가중 평균이고, 평균은 가장 작은 값보다 작아질 수 없습니다. 값이 0.6, 0.8, 1.0인 세 점이면, 무게를 1, 0, 0으로 주면 0.6, 반반씩이면 0.7, 3분의 1씩이면 0.8입니다. 반 평균이 꼴찌 점수보다 낮을 수 없는 것과 같습니다. 그래서 가장 작게 하려면 가장 가까운 점에 무게를 다 몰아줍니다. w가 정해져 있으면 람다는 argmin 점 하나에 1입니다. 그런데 답에서는 서포트 벡터 세 점이 똑같이 가장 가까워서, 0.67, 0.67, 0.67처럼 동점입니다. 이때는 셋에 어떻게 나눠 줘도 같은 값이라 무게를 나눠 가질 수 있습니다. 나머지 17개는 더 멀어서 무게가 0입니다. 나누는 비율 0.5, 0.25, 0.25는 학습한 w 안의 몫과 같습니다. 이어서 확인합니다.}
\end{frame}

% --- argmin 이 여럿: 서포트 벡터 몇 개로 쓸 수 있다
\begin{frame}{가장 가까운 점이 여럿: $\mathbf{w}$ 는 서포트 벡터 몇 개로 쓸 수 있다}
\vfill
\centering
{\large 한 걸음 \;=\; 가장 가까운 점(argmin)의 $y_k \mathbf{x}_k$ 를 \textbf{더한다} \;$\to$\; $\mathbf{w}$ 는 점들을 더한 것}\\[0.4em]
{\large 답 근처: \textbf{세 점이 똑같이 가장 가깝다} (값 0.67) \;$\to$\; 번갈아 argmin 이 되어 \textbf{번갈아 민다}}\\[0.8em]
\renewcommand{\arraystretch}{1.35}
{\normalsize\begin{tabular}{r|ccc|c}
 & \textcolor{ksared}{빨강 $r$} & \textcolor{ksablue}{파랑 $b_1$} & \textcolor{ksablue}{파랑 $b_2$} & 나머지 17개\\ \hline
뒤 1만 걸음 동안 argmin 이 된 횟수 & 5025 & 2511 & 2512 & 0\\
$\mathbf{w}$ 안의 몫 (2만 걸음) & \kb{0.5} & \kb{0.25} & \kb{0.25} & 0.000002
\end{tabular}}\\[0.8em]
{\large $\mathbf{w} \propto 0.5\, r - 0.25\, b_1 - 0.25\, b_2 = (0.475,\ 0.475)$ \;$\to$\; 길이 $1$ 로 \; $(0.71,\ 0.71)$ \; $=$ 학습한 답}\\[0.4em]
{\large \kq{서포트 벡터 3개만으로} $\mathbf{w}$ 를 쓸 수 있다 \;(나머지는 지워도 같은 선)}
\vfill
\note{학습한 w가 무엇으로 이루어져 있는지 봅시다. 한 걸음마다 한 일은 가장 가까운 점, argmin 점의 y x를 더하고 길이를 1로 맞춘 것뿐입니다. 그래서 w는 점들을 더해 만든 것입니다. 답 근처에서는 세 점, 빨강 r과 파랑 b1, b2가 똑같이 0.67로 가장 가깝습니다. 그래서 빨강을 밀면 파랑이 가장 가까워지고, 파랑을 밀면 빨강이 가까워지면서 번갈아 argmin이 됩니다. 마지막 1만 걸음 동안 빨강은 5025번, 파랑 두 점은 2511번, 2512번 argmin이었고 나머지 17개는 한 번도 아니었습니다. 그래서 w 안에 들어 있는 몫도 0.5 대 0.25 대 0.25이고, 나머지 17개의 몫은 처음 평균 선에서 받은 것이 나눌 때마다 줄어 0.000002, 사실상 0입니다. 실제로 0.5 r 빼기 0.25 b1 빼기 0.25 b2는 0.475, 0.475이고, 길이를 1로 맞추면 0.71, 0.71로 학습한 답과 같습니다. 빨강이 0.5인 것은 빨강 한 점이 파랑 두 점 몫을 혼자 버텨야 선이 한쪽으로 기울지 않기 때문입니다. 그래서 w는 서포트 벡터 세 점만으로 쓸 수 있고, 나머지 점은 지워도 같은 선이 나옵니다. 다음 시간에는 이 몫을 직접 다루는 방법을 봅니다.}
\end{frame}

\section{크기 약속}
% --- 크기 문제: 같은 선, 숫자는 무한히 많다
\begin{frame}{같은 선을 쓰는 숫자는 무한히 많다}
\vfill
\centering
{\large $\displaystyle \max_{\mathbf{w},\,b,\ \lVert\mathbf{w}\rVert = 1}\ \min_{k}\ y_k D(\mathbf{x}_k)$ \qquad 왜 $\lVert\mathbf{w}\rVert = 1$ 로 묶었나?}\\[1.1em]
\renewcommand{\arraystretch}{1.45}
{\normalsize\begin{tabular}{l|c|c|c|c}
$\mathbf{w}$, \; $b$ & $\lVert\mathbf{w}\rVert$ & $D(\mathbf{x}_r)$ & 거리 $= D(\mathbf{x}_r)/\lVert\mathbf{w}\rVert$ & 묶은 것\\ \hline
$(0.71,\ 0.71)$, \; $0.04$ & 1 & 0.67 & \kb{0.67} & $\lVert\mathbf{w}\rVert = 1$ (지금까지)\\
$(1.41,\ 1.41)$, \; $0.07$ \;($\times 2$) & 2 & 1.34 & \kb{0.67} & 없음\\
$(1.05,\ 1.05)$, \; $0.05$ & 1.49 & \kq{1} & \kb{0.67} & $D(\mathbf{x}_r) = 1$ (새 약속)
\end{tabular}}\\[1em]
{\large 선도 거리도 \textbf{그대로}, 숫자 크기만 다르다 $\to$ 하나를 고를 \kq{약속}이 필요하다}\\[0.3em]
{\small\color{gray} $\mathbf{x}_r$ = \textbf{가장 가까운} 빨간 점 서포트 벡터 $(0.7,\ 0.2)$}
\vfill
\note{방금까지의 식은 w의 길이를 1로 묶고, 가장 가까운 점의 거리를 최대로 하는 것이었습니다. 왜 길이를 묶었을까요. 같은 선을 쓰는 방법이 무한히 많기 때문입니다. w와 b에 같은 수를 곱해도 D x가 0인 곳은 그대로라 선이 같습니다. 표를 봅시다. 첫 줄은 방금 학습한 답입니다. 길이가 1이고 빨간 서포트 벡터의 값이 0.67, 거리도 0.67입니다. 둘째 줄은 두 배를 한 것입니다. 길이는 2, 값은 1.34로 두 배가 되지만, 거리는 값을 길이로 나눈 것이라 여전히 0.67입니다. 셋째 줄은 서포트 벡터의 값이 딱 1이 되게 크기를 고른 것입니다. 길이는 1.49이고 거리는 역시 0.67입니다. 세 줄 모두 같은 선입니다. 그래서 무한히 많은 것 중 하나를 고르는 약속이 필요합니다. 지금까지의 약속은 w의 길이가 1이었습니다. 원 논문은 여기서 약속을 바꿉니다.}
\end{frame}

% --- M‖w‖ = 1 → M = 1/‖w‖
\begin{frame}{약속을 $M\lVert\mathbf{w}\rVert = 1$ 로 바꾸면: $M = 1/\lVert\mathbf{w}\rVert$}
\vfill
\centering
\renewcommand{\arraystretch}{1.0}
{\large\begin{tabular}{rl}
\textbf{새 약속} & $M\lVert\mathbf{w}\rVert = 1$ \quad $\to$ \quad \fcolorbox{fbeta}{fbeta!10}{$\displaystyle M = \frac{1}{\lVert\mathbf{w}\rVert}$}\\[1.05em]
\textbf{조건} & 모든 점의 거리가 $M$ 이상: \; $\dfrac{y_k D(\mathbf{x}_k)}{\lVert\mathbf{w}\rVert} \ge M = \dfrac{1}{\lVert\mathbf{w}\rVert}$ \; $\to$ \; \kb{$y_k D(\mathbf{x}_k) \ge 1$}\\[1.05em]
\textbf{목표} & $M$ 최대 \; $=$ \; $\dfrac{1}{\lVert\mathbf{w}\rVert}$ 최대 \; $=$ \; \kb{$\lVert\mathbf{w}\rVert$ 최소}
\end{tabular}}\\[0.8em]
\fcolorbox{fbeta}{fbeta!8}{\large $\displaystyle \min_{\mathbf{w},\,b}\ \tfrac{1}{2}\lVert\mathbf{w}\rVert^2$ \quad 조건: $y_k D(\mathbf{x}_k) \ge 1$, \; $k = 1, \dots, p$}\\[0.35em]
{\small\color{gray} $\tfrac12$ 과 제곱은 미분하기 쉬우라고 (답은 같다)}\\[0.6em]
{\large 풀고 나면 마진 \; $M^* = \dfrac{1}{\lVert\mathbf{w}^*\rVert} = \dfrac{1}{1.49} = \kb{0.67}$}
\vfill
\note{원 논문은 w의 길이를 1로 묶는 대신, 마진 M과 w의 길이의 곱을 1로 묶습니다. 양변을 w의 길이로 나누면 M은 1 나누기 w의 길이입니다. 이것이 핵심입니다. 마진이 w의 길이 하나로 정해집니다. 이제 조건을 다시 씁니다. 모든 점은 경계에서 M 이상 떨어져야 합니다. 점의 거리는 y k D x k 나누기 w의 길이이고, M 자리에 1 나누기 w의 길이를 넣으면 양쪽에 같은 분모가 있어서 지워집니다. 남는 것은 y k D x k가 1 이상입니다. 목표도 다시 씁니다. M을 최대로 하는 것은 1 나누기 w의 길이를 최대로 하는 것이고, 분수를 크게 하려면 분모를 작게 해야 하니, w의 길이를 최소로 하는 것입니다. 그래서 아래 상자의 문제가 됩니다. 원 논문은 w 길이 제곱을 최소로 한다고 쓰고, 뒤의 라그랑주 식에서 2분의 1을 붙입니다. 길이가 가장 작은 w는 제곱도 가장 작고 2분의 1을 곱해도 마찬가지라 답은 같습니다. 이 문제를 풀어 w 별을 구하면 마진은 1 나누기 w 별의 길이입니다. 이 데이터에서는 1 나누기 1.49, 0.67입니다. 앞에서 학습한 것과 똑같은 마진입니다. 다음 시간은 이 식에서 출발합니다.}
\end{frame}

\section{정리}
% --- 다중 클래스
\begin{frame}{클래스가 3개 이상이면}
\vfill
\centering
\renewcommand{\arraystretch}{1.6}
{\large\begin{tabular}{c|c|c}
 & {\Large\kb{일대일 (one-vs-one)}} & {\Large\kq{일대다 (one-vs-all)}}\\ \hline
만드는 법 & 클래스 두 개씩 짝마다 SVM & 클래스마다 ``나 vs 나머지'' SVM\\
개수 & $c(c-1)/2$ 개 & $c$ 개\\
예측 & \textbf{다수결} & $D$ 가 가장 큰 클래스\\ \hline
숫자 0$\sim$9 ($c = 10$) & \kb{45개} & \kq{10개}
\end{tabular}}
\vfill
\note{SVM은 초평면 하나로 두 쪽을 가르는 방법이라 기본은 두 클래스용입니다. 클래스가 c개면 두 가지 방법을 씁니다. 일대일은 클래스 두 개씩 짝을 지어 짝마다 SVM을 하나씩 만듭니다. c 곱하기 c 빼기 1 나누기 2개가 필요하고, 새 점은 모든 짝에서 이긴 클래스를 세어 가장 많이 이긴 클래스로 정합니다. sklearn의 SVC가 안에서 이 방식을 씁니다. 일대다는 클래스마다 나 대 나머지 전부로 SVM을 하나씩, 모두 c개 만들고, 새 점에서 D 값이 가장 큰, 즉 경계에서 가장 멀리 확신하는 클래스로 정합니다. 나머지 쪽이 훨씬 많아 불균형이 생길 수 있습니다. 손글씨 숫자 0부터 9까지면 일대일은 45개, 일대다는 10개입니다. SVM 원 논문의 손글씨 숫자 실험도 숫자마다 하나씩, 경계 10개를 만든 일대다였습니다.}
\end{frame}

\section{연습}
% --- 실습
\begin{frame}{실습}
\vfill
\centering\begin{tabular}{rl}
\large 1 & 슬라이드의 점 20개로 선형 SVM: $M$ 0.67 (폭 1.34), 서포트 벡터 3개, 먼 점 옮기기\\[0.5em]
\large 2 & 슬라이드의 학습 방법을 직접 짜서 sklearn 과 같은 선인지 확인
\end{tabular}\vskip1em
\href{https://colab.research.google.com/github/SuminHan/book-ml/blob/main/slides/ml1/week02/2_svm/lab/svm.ipynb}{\textcolor{ksagolddk}{▶ Colab에서 실습 노트북 열기}}
\vfill
\note{실습은 두 부분입니다. 첫째, 슬라이드의 z-score 한 점 20개로 선형 SVM을 학습해서 마진 0.67, 길 폭 1.34와 서포트 벡터 3개를 재현하고, 먼 점을 옮겨도 경계가 그대로인지 확인합니다. 둘째, 슬라이드의 학습 방법, 가장 가까운 점이 밀고 w 길이를 1로 맞추는 것을 직접 코드로 짜서 sklearn과 같은 선이 나오는지 비교합니다. 소프트 마진은 다음 시간, 커널과 지폐 데이터는 그다음 시간 실습에서 합니다.}
\end{frame}

% --- 정리
\begin{frame}{정리}
\vfill
\centering\normalsize\begin{tabular}{l}
\textbf{초평면} $D(\mathbf{x}) = \mathbf{w}\cdot\mathbf{x} + b = 0$: 값의 부호로 양쪽을 나눈다\\[0.5em]
\textbf{최대 마진}: $M^* = \max_{\lVert\mathbf{w}\rVert = 1} \min_k y_k D(\mathbf{x}_k)$, 경계는 \kb{서포트 벡터}가 정한다\\[0.5em]
\textbf{학습}: z-score 한 뒤, 가장 가까운 점이 밀고 길이 1 로 맞추기를 반복\\[0.5em]
\textbf{다중 클래스}: 일대일 $c(c-1)/2$ 개 · 일대다 $c$ 개\\[0.5em]
\textbf{새 약속}: $M\lVert\mathbf{w}\rVert = 1$ $\to$ $\min \tfrac12\lVert\mathbf{w}\rVert^2$, 조건 $y_k D(\mathbf{x}_k) \ge 1$\\[0.5em]
\kq{다음}: 점이 겹치면? $\to$ 소프트 마진 \quad 직선으로 못 가르면? $\to$ 커널
\end{tabular}
\vfill
\note{정리합니다. 초평면은 값의 부호로 공간을 둘로 나눕니다. 최대 마진 분류기는 가장 가까운 점까지의 거리 M을 최대로 하는 초평면이고, w의 길이를 1로 맞추면 값이 곧 거리입니다. 원 논문의 한 줄로는 M 별이 max min y k D x k입니다. 경계는 서포트 벡터 몇 개가 정합니다. 학습은 점들을 z-score 한 뒤 매 걸음 가장 가까운 점이 선을 밀고 w의 길이를 다시 1로 맞추는 것을 반복하고, sklearn은 SMO라는 다른 풀이법으로 같은 답을 냅니다. 클래스가 여러 개면 일대일이나 일대다로 SVM을 여러 개 만듭니다. 그리고 같은 선을 쓰는 약속을 M 곱하기 w의 길이가 1로 바꾸면, 마진 최대는 w의 길이 최소가 됩니다. 오늘 데이터는 직선으로 깔끔하게 갈렸습니다. 다음 시간에는 두 색이 겹치면 어떻게 하는지 소프트 마진을, 그다음 시간에는 직선으로 못 가르는 데이터를 다루는 커널을 배웁니다.}
\end{frame}
% --- 참고 문헌
\begin{frame}{참고 문헌}
\vfill
\small
\begin{itemize}
\setlength{\itemsep}{0.7em}
\item Boser, B. E., Guyon, I. M., \& Vapnik, V. N. (1992). A training algorithm for optimal margin classifiers. \textit{COLT '92}, 144--152. \paperlink{https://doi.org/10.1145/130385.130401}
\item Cortes, C., \& Vapnik, V. (1995). Support-vector networks. \textit{Machine Learning}, 20(3), 273--297. \paperlink{https://doi.org/10.1007/BF00994018}
\item James, G., Witten, D., Hastie, T., \& Tibshirani, R. \textit{An Introduction to Statistical Learning}, 9장. \paperlink{https://doi.org/10.1007/978-1-0716-1418-1}
\item Wang, S. (2019). CME 250 Lecture 5: Support Vector Machines, Stanford. \paperlink{https://web.stanford.edu/class/cme250/files/cme250_lecture5.pdf}
\end{itemize}
\vfill
\note{이 덱의 참고 문헌입니다.}
\end{frame}

\end{document}
